% Libro 4 · Mundo 1 · Inferencia (app: pu-inferencia)
\mundo{Inferencia}{pu-inferencia}{Inferir es leer entre líneas: sacar de lo que el texto dice algo que no dice pero que
se sigue con seguridad. No es adivinar ni opinar. La diferencia entre una buena inferencia y una especulación es que la
primera tiene \emph{pruebas} en el texto. Este es el nivel que más pesa en la prueba.}{%
\item \textbf{1.1} Ideas implícitas
\item \textbf{1.2} Presupuestos: lo que se da por sentado
\item \textbf{1.3} Causa y efecto
\item \textbf{1.4} Actitudes y motivos de los personajes}

% ------------------------------------------------------------------
\leccion{pu-in-1}{Ideas implícitas}{Juego de encontrar la pista que sostiene una conclusión.}

\begin{reto}
«Cuando Laura llegó a la estación, el andén estaba vacío y el reloj marcaba las 7:10. Se sentó en una banca y sacó un libro
largo.» ¿Qué se puede inferir?
\opcionesl{Laura perdió el tren y esperará el siguiente.}{Laura trabaja en la estación.}{El tren llegó temprano.}{Laura odia leer.}
\end{reto}

\tecnica{La técnica: conclusión con pruebas}
Una \clave{inferencia válida} cumple dos condiciones: (1) no está dicha literalmente, y (2) se sigue de lo dicho sin
necesidad de inventar. Para comprobarla, pregúntate: \emph{¿qué frase del texto me obliga a pensar esto?} Si no encuentras
ninguna, es especulación.

\figura{pistas}{Una inferencia se apoya en pistas que sí están en el texto.}

\subsubsection{Grados de inferencia}
Entre «el texto lo dice» y «el texto no dice nada de eso» hay una zona: lo que el texto \emph{sugiere con fuerza}. El
examen premia la opción que está más cerca de lo seguro. Una opción con «probablemente», «sugiere», «puede» suele ser más
defendible que una con «siempre» o «es seguro que».

\begin{regla}
\begin{itemize}
  \item Señala la frase que prueba la inferencia. Sin prueba, no hay inferencia.
  \item Prefiere la conclusión mínima: la que exige suponer menos.
  \item Cuidado con lo que «suena lógico» por conocimiento propio, pero no tiene apoyo en el texto.
\end{itemize}
\end{regla}

\begin{texto}{La vecina del cuarto piso}{Relato · texto original escrito para este libro}
\pa Doña Inés vive sola en el cuarto piso desde que yo tengo memoria. Todos los días, a las cinco de la tarde, baja a la
portería con una bolsa de pan y se queda conversando con el celador hasta que oscurece. Él dice que es la única hora en que
ella habla con alguien.

\pa El mes pasado no bajó durante una semana. El celador subió dos veces a tocarle la puerta; la segunda vez ella abrió en
bata, con la voz ronca, y le dijo que no se preocupara, que era solo una gripa. Desde entonces, él le sube el pan.

\pa Ayer la vi de nuevo en la portería. Llevaba la bolsa de pan, pero esta vez traía dos, y una se la entregó al celador
sin decir nada. Él la recibió también en silencio, y los dos se quedaron mirando la calle hasta que encendieron los
faroles.
\end{texto}

\subsection{Preguntas comentadas}
\begin{comentada}{¿Qué se puede inferir del gesto de doña Inés en el \S\,3?}
\mal{A}{Que compró pan de más por error.}{No hay ninguna prueba; contradice la intención del gesto.}
\bien{B}{Que agradece al celador lo que hizo por ella cuando estuvo enferma.}{El \S\,2 cuenta que él le subía el pan; ahora
ella le trae una bolsa: el gesto devuelve el favor.}
\mal{C}{Que el celador no tiene dinero para comer.}{\emph{No lo dice} ni se sigue.}
\mal{D}{Que ella y el celador están peleados.}{\emph{Dice lo contrario}: comparten la tarde.}
\end{comentada}

\begin{trampa}
Elegir una inferencia emocionante pero sin pruebas («están enamorados»). Puede que sí, pero el texto no lo sostiene; la
gratitud sí.
\end{trampa}

\taller
\piso{Como en el examen}{opción múltiple}
\begin{ejercicios}
\item Del \S\,1 se puede inferir que doña Inés
\opcionesl{tiene muchos amigos en el edificio.}{probablemente se siente sola.}{trabaja de noche.}{no le gusta el pan.}
\sol{B. «la única hora en que ella habla con alguien».}
\item El silencio con que se intercambian el pan en el \S\,3 sugiere que
\opcionesl{están enojados.}{entre ellos hay una confianza que no necesita palabras.}{no se conocen.}{el celador no quería el pan.}
\sol{B.}
\item ¿Cuál inferencia NO tiene apoyo en el texto?
\opcionesl{Doña Inés estuvo enferma.}{El celador se preocupó por ella.}{Doña Inés tiene hijos que la visitan.}{Doña Inés recuperó su rutina.}
\sol{C.}
\end{ejercicios}

\begin{retoresuelto}
Andén vacío (el tren ya pasó o no ha llegado), y Laura se prepara para una espera larga (un libro largo): lo más defendible
es que perdió el tren y esperará el siguiente. Las demás no tienen apoyo. \textbf{Opción A.}
\end{retoresuelto}
\bitacora

% ------------------------------------------------------------------
\leccion{pu-in-2}{Presupuestos}{Juego de descubrir lo que una frase da por sentado.}

\begin{reto}
«Andrés dejó de fumar hace un año.» Esta frase presupone que
\opcionesd{Andrés fumaba antes.}{Andrés está enfermo.}{fumar es bueno.}{Andrés nunca fumó.}
\end{reto}

\tecnica{La técnica: lo que se da por hecho}
Un \clave{presupuesto} es una información que la frase no afirma, pero que da por sentada para tener sentido. «Dejó de
fumar» presupone que fumaba; «otra vez llegó tarde» presupone que ya había llegado tarde; «lamento que no vinieras»
presupone que no viniste. Una prueba sencilla: si se niega la frase («no dejó de fumar»), el presupuesto sigue en pie
(fumaba antes, sigue fumando).

\subsubsection{Marcas de presupuestos}
Verbos de cambio (\emph{dejar de, empezar a, seguir}), adverbios como \emph{otra vez, todavía, ya, también}, verbos de
sentimiento o conocimiento (\emph{lamentar, saber, reconocer que}), y descripciones como \emph{el nuevo director}
(presupone que hubo otro).

\begin{regla}
\begin{itemize}
  \item Busca verbos de cambio, «otra vez», «todavía», «también», «el nuevo…».
  \item Prueba de la negación: el presupuesto sobrevive cuando se niega la frase.
  \item En argumentos, los presupuestos son ideas que el autor no defiende porque cree que todos las aceptan.
\end{itemize}
\end{regla}

\begin{texto}{La nueva biblioteca del colegio}{Carta de un estudiante · texto original escrito para este libro}
\pa Señora rectora: le escribo en nombre del consejo estudiantil porque, \emph{otra vez}, la biblioteca permanecerá cerrada
durante las vacaciones de mitad de año. Sabemos que la nueva bibliotecaria todavía se está adaptando, pero muchos
estudiantes necesitamos ese espacio precisamente en vacaciones, cuando preparamos el examen de admisión.

\pa Lamentamos que la decisión se haya tomado sin consultarnos. Proponemos que la biblioteca abra al menos tres mañanas
por semana, con estudiantes voluntarios del último grado que ayuden a cuidarla.

\pa También le pedimos que considere ampliar la sección de libros de preparación, que sigue siendo la más pequeña de la
biblioteca. Quedamos atentos a su respuesta.
\end{texto}

\subsection{Preguntas comentadas}
\begin{comentada}{La expresión «otra vez» \pf{1} presupone que}
\mal{A}{la rectora está enojada.}{\emph{No lo dice}.}
\bien{B}{la biblioteca ya había estado cerrada en vacaciones antes.}{«Otra vez» indica repetición.}
\mal{C}{la biblioteca nunca abre.}{\emph{Dice de más}.}
\mal{D}{los estudiantes no estudian en vacaciones.}{\emph{Dice lo contrario}.}
\end{comentada}

\taller
\piso{Como en el examen}{opción múltiple}
\begin{ejercicios}
\item «La nueva bibliotecaria» presupone que
\opcionesl{antes hubo otra bibliotecaria.}{la bibliotecaria es joven.}{la biblioteca es nueva.}{la bibliotecaria no sabe su oficio.}
\sol{A.}
\item «Lamentamos que la decisión se haya tomado sin consultarnos» presupone que
\opcionesl{los estudiantes fueron consultados.}{la decisión se tomó sin consultar a los estudiantes.}{la rectora consultará en el futuro.}{la decisión es ilegal.}
\sol{B.}
\item «que sigue siendo la más pequeña» \pf{3} presupone que
\opcionesl{la sección ya era la más pequeña antes.}{la sección creció.}{no hay libros de preparación.}{la biblioteca es pequeña.}
\sol{A.}
\end{ejercicios}

\begin{retoresuelto}
«Dejar de» presupone que antes se hacía. Aunque se niegue («no dejó de fumar»), sigue presupuesto que fumaba.
\textbf{Opción A.}
\end{retoresuelto}
\bitacora

% ------------------------------------------------------------------
\leccion{pu-in-3}{Causa y efecto}{Juego de las fichas que se encadenan.}

\begin{reto}
«Las lluvias de marzo dañaron las vías, por lo que los camiones tardaron el doble en llegar al mercado y los precios de
la papa subieron.» Según el fragmento, ¿cuál fue la causa directa del alza de precios?
\opcionesd{las lluvias}{la demora de los camiones}{el mercado}{la papa}
\end{reto}

\tecnica{La técnica: dibujar la cadena}
Los textos encadenan hechos: A produce B, que produce C. Las preguntas piden a veces la causa \emph{directa} (la
inmediatamente anterior) y a veces la causa \emph{original}. Dibujar la cadena evita confundirlas:

\figura{cadena}{Una cadena causal: cada hecho empuja al siguiente.}

Además, hay que distinguir la causa de la simple \emph{coincidencia} (lo que ocurrió al mismo tiempo) y de la
\emph{condición} (lo que hizo posible, pero no provocó).

\begin{regla}
\begin{itemize}
  \item Marcas de causa: porque, ya que, debido a, gracias a, a causa de. De efecto: por eso, así que, de modo que,
  provocó, generó.
  \item Dibuja la cadena y fíjate si piden la causa directa o la original.
  \item Ocurrir antes no es causar.
\end{itemize}
\end{regla}

\begin{texto}{El año sin luciérnagas}{Divulgación · texto original escrito para este libro}
\pa En las veredas de algunos municipios, los mayores recuerdan noches llenas de luciérnagas. Hoy, muchos niños de esas
mismas veredas nunca han visto una.

\pa Las luciérnagas usan su luz para encontrar pareja: los machos vuelan encendiéndose y apagándose, y las hembras
responden desde el suelo. Cuando la noche está iluminada por bombillos y faroles, esas señales se pierden, y menos parejas
logran encontrarse. Con menos encuentros, nacen menos larvas.

\pa Las larvas, además, viven en el suelo húmedo y se alimentan de caracoles y babosas. Por eso, cuando se secan los
humedales o se usan venenos contra caracoles en los cultivos, también se quedan sin hogar y sin alimento.

\pa La buena noticia es que las causas son conocidas. Apagar las luces innecesarias, dirigir los faroles hacia el suelo y
conservar zonas húmedas son medidas sencillas que, en algunos lugares, ya han devuelto las luces de las noches.
\end{texto}

\subsection{Preguntas comentadas}
\begin{comentada}{Según el \S\,2, ¿por qué nacen menos larvas cuando la noche está iluminada?}
\mal{A}{Porque la luz mata a las larvas.}{\emph{No lo dice}.}
\bien{B}{Porque las señales de luz se pierden y menos parejas se encuentran.}{Cadena: luz artificial $\to$ señales
perdidas $\to$ menos encuentros $\to$ menos larvas.}
\mal{C}{Porque las hembras vuelan encendiéndose.}{\emph{Dice otra cosa}: son los machos los que vuelan.}
\mal{D}{Porque no hay caracoles.}{Es otra causa, del \S\,3, no la de la luz.}
\end{comentada}

\taller
\piso{Como en el examen}{opción múltiple}
\begin{ejercicios}
\item Según el \S\,3, el uso de venenos contra caracoles afecta a las luciérnagas porque
\opcionesl{envenena a los adultos.}{deja a las larvas sin alimento.}{ilumina los cultivos.}{seca los humedales.}
\sol{B.}
\item ¿Cuál de las siguientes NO aparece como causa de la disminución de luciérnagas?
\opcionesd{la luz artificial}{la desecación de humedales}{el cambio de temperatura}{los venenos para caracoles}
\sol{C.}
\item El \S\,4 presenta las medidas como
\opcionesl{costosas e imposibles.}{sencillas y, en algunos lugares, ya efectivas.}{obligatorias por ley.}{inútiles.}
\sol{B.}
\end{ejercicios}

\begin{retoresuelto}
Cadena: lluvias $\to$ vías dañadas $\to$ camiones tardan el doble $\to$ precios suben. La causa \emph{directa} del alza es
la demora; las lluvias son la causa original. \textbf{Opción B.}
\end{retoresuelto}
\bitacora

% ------------------------------------------------------------------
\leccion{pu-in-4}{Actitudes y motivos}{Juego de leer las emociones de un personaje por lo que hace.}

\begin{reto}
«Camilo leyó la carta dos veces, la dobló con cuidado y la guardó en el bolsillo de la camisa, del lado del corazón.» La
actitud de Camilo hacia la carta es de
\opcionesd{indiferencia}{aprecio}{enojo}{burla}
\end{reto}

\tecnica{La técnica: los hechos delatan}
Los buenos textos rara vez dicen «estaba triste»: lo muestran. Para inferir una \clave{actitud} o un \clave{motivo}, se
miran las acciones, los gestos, las palabras elegidas y los silencios. Luego se busca la opción que nombre esa emoción
con la intensidad justa: ni más (furia, si solo hay molestia) ni menos.

\begin{regla}
\begin{itemize}
  \item Haz una lista de las acciones y palabras del personaje.
  \item Pregunta: ¿qué sentimiento explica \emph{todas} esas acciones a la vez?
  \item Ajusta la intensidad: preocupación no es pánico; molestia no es odio.
\end{itemize}
\end{regla}

\begin{texto}{El primer día}{Relato · texto original escrito para este libro}
\pa Mi papá me acompañó hasta la puerta de la universidad aunque le dije tres veces que no hacía falta. En el bus no habló
mucho; miraba por la ventana y se arreglaba el cuello de la camisa, que no estaba torcido.

\pa Cuando llegamos, se quedó en la acera. Me dio la mano como a un adulto, pero no me la soltó enseguida. «Usted es la
primera de la familia que llega aquí», dijo, y después tosió como si se le hubiera metido algo en la garganta.

\pa Caminé hacia el edificio sin voltear. Al llegar a la escalera miré hacia atrás: seguía en la acera, con las manos en
los bolsillos, sin moverse. Esa tarde, cuando volví a la casa, mi mamá me contó que él había llegado al trabajo dos horas
tarde.
\end{texto}

\subsection{Preguntas comentadas}
\begin{comentada}{¿Qué actitud muestra el padre en el texto?}
\mal{A}{Fastidio por tener que acompañarla.}{\emph{Dice lo contrario}: él insistió en ir.}
\bien{B}{Orgullo y emoción contenida.}{Se arregla la camisa (nervios), no suelta la mano, dice la frase sobre «la
primera de la familia», tose para disimular, se queda inmóvil y llega tarde al trabajo.}
\mal{C}{Indiferencia.}{Todas sus acciones la contradicen.}
\mal{D}{Miedo a que le pase algo malo.}{\emph{No lo dice}: nada sugiere peligro.}
\end{comentada}

\taller
\piso{Como en el examen}{opción múltiple}
\begin{ejercicios}
\item La tos del padre \pf{2} sugiere que
\opcionesl{está enfermo.}{intenta disimular la emoción.}{tiene frío.}{está molesto.}
\sol{B.}
\item El dato final, que llegó dos horas tarde al trabajo, permite inferir que
\opcionesl{se perdió en el camino.}{se quedó un buen rato frente a la universidad.}{no tenía trabajo ese día.}{el bus se dañó.}
\sol{B. Coherente con «seguía en la acera, sin moverse».}
\item ¿Por qué la narradora caminó «sin voltear» \pf{3}?
\opcionesl{Porque estaba enojada con su padre.}{Probablemente para mostrarse fuerte o no emocionarse.}{Porque no lo vio.}{Porque llegaba tarde.}
\sol{B. Luego, en la escalera, sí miró atrás.}
\end{ejercicios}

\begin{retoresuelto}
Leerla dos veces, doblarla con cuidado y guardarla «del lado del corazón» muestra que la valora. \textbf{Opción B.}
\end{retoresuelto}
\bitacora

% ------------------------------------------------------------------
\evaluacion
\begin{texto}{La última función}{Crónica · texto original escrito para este libro (el teatro es imaginario)}
\pa El teatro Colón de San Pedro cerró sus puertas un domingo de lluvia. No hubo discursos: el dueño, don Alirio, apagó
las luces de la marquesina a las diez de la noche, como siempre, y se quedó un rato mirando la fachada desde la otra
acera.

\pa Durante cincuenta años, el Colón fue el único cine del pueblo. Allí se dieron los primeros besos de varias
generaciones, allí se lloró con las películas mexicanas y se gritó con las de vaqueros. En los años noventa todavía se
llenaba los sábados; en los últimos tiempos, a veces la función empezaba con tres espectadores.

\pa Don Alirio nunca subió el precio de la boleta, aunque sus hijos se lo pedían. Decía que el cine era «para que la gente
saliera de la casa», no para hacerse rico. Cuando las plataformas de películas llegaron a los celulares, la gente dejó de
salir.

\pa La última película fue una comedia. Los pocos asistentes se rieron en los mismos momentos, y al final, cuando se
encendieron las luces, nadie se levantó enseguida.
\end{texto}
\begin{ejercicios}
\item Del \S\,1 se puede inferir que don Alirio
\opcionesl{estaba feliz de cerrar.}{se despidió del teatro con tristeza contenida.}{no quería volver nunca.}{olvidó apagar las luces.}
\sol{B. (1.4)}
\item Según el texto, la causa principal de que la gente dejara de ir al cine fue
\opcionesl{el precio de la boleta.}{la llegada de las películas a los celulares.}{la lluvia.}{las películas mexicanas.}
\sol{B. \S\,3; el precio nunca subió. (1.3)}
\item La frase «el cine era para que la gente saliera de la casa» presupone que
\opcionesl{la gente salía poco de la casa.}{el cine tenía una función más social que comercial para don Alirio.}{don Alirio era rico.}{las películas eran malas.}
\sol{B. (1.2)}
\item Que nadie se levantara enseguida al final \pf{4} sugiere que
\opcionesl{la película fue aburrida.}{los asistentes querían prolongar la despedida del teatro.}{estaban dormidos.}{esperaban otra película.}
\sol{B. (1.1)}
\item ¿Cuál inferencia NO se puede sostener?
\opcionesl{El teatro fue importante en la vida del pueblo.}{Don Alirio priorizaba su idea del cine sobre la ganancia.}{Los hijos de don Alirio querían subir el precio.}{El teatro será demolido.}
\sol{D. (1.1)}
\end{ejercicios}
